\newpage
\hypertarget{para:amorun}{}

\mypoemtitle{Amor è un disio che ven da core - Giacomo da Lentini}{Sonetto}

\textbf{Commento introduttivo:}\\
In questo sonetto Giacomo da Lentini applica una fenomenologia dell'amore sull'uomo descrivendo quello che accade quando un'uomo si innamora ed andando ad avvalorare l'importanza della vista in questo processo

\vspace{1.5em}

\begin{adjustwidth}{2em}{0pt}
\poemlines{5}
\verselinenumbersleft
\begin{verse}
Amore è uno desi[o] che ven da’ core \\
per abondanza di gran piacimento; \\
e li occhi in prima genera[n] l’amore \\
e lo core li dà nutricamento. \\!

Ben è alcuna fiata om amatore \label{v5}\\
senza vedere so ’namoramento, \\
ma quell’amor che stringe con furore \\ 
da la vista de li occhi ha nas[ci]mento: \\!

ché li occhi rapresenta[n] a lo core \\
d’onni cosa che veden bono e rio  \\
com’è formata natural[e]mente; \label{v11}\\!

e lo cor, che di zo è concepitore, \\
imagina, e [li] piace quel desio: \\
e questo amore regna fra la gente. \label{v14}\\
\end{verse}
\end{adjustwidth}

\vspace{1.5em}

\begin{commentaries}
    \scolio[\ref{v5}]{fiata} Provenzalismo: volta
    \scolio[\ref{v11}]{naturalmente} L'amore è un'esperienza fisica naturale
    \scolio[\ref{v14}]{regna fra la gente} è diffuso in tutto il mondo
\end{commentaries}

\vspace{0.5em}
{\color{Tomato!40}\hrule height 0.5pt}
\vspace{1em}

\subsubsection*{Analisi}
\begin{quote}
\small 
Il componimento si presenta come risposta ad una \textbf{tenzone} cominciata da Jacopo Mostacchi. A questa tenzone prendono parte Jacopo mostacchi, Giacomo da Lentini e Pier delle Vigne i quali cercano di spiegare in maniera filosofica quale sia la natura dell'amore. In particolare fanno riferimento ai concetti aristotelici di \textbf{sostanza} ed \textbf{accidente}. Per Aristotele la sostanza di un oggetto è quella che lo rende tale, non può quindi esistere tale oggetto senza la sua sostanza, un accidente è invece un'attributo di questo elemento che può variare e pur identificandolo non lo rende o meno tale.
Jacopo Mostacchi sostiene che l'amore sia un accidente, Pier delle Vigne sostiene che sia una sostanza senza cui l'uomo non potrebbe vivere e Giacomo da Lentini si pone un po' a metà tra le due visioni.

La struttura di questo sonetto è abbastanza classica, troviamo infatit nella prima quartina la tesi, nella seconda un'antitesi e nelle terzine degli argomenti volti a dare valore alla tesi.

Già a partire dal vv.1 possiamo notare la natura didascalica di questo componimento: Giacomo sembra spiegare al lettore cosa sia l'amore e come funzioni mettendo in luce due componenti fondamentali: \textbf{il cuore} in cui il sentimento risiede e da cui viene nutrito e \textbf{gli occhi} che hanno l'importante compito di generare il sentimento amoroso, tema già affrontato in altri testi dei Siciliani.
Ritornando a considerare la nostra tenzone potremmo sostenere che Giacomo sia più per definire l'amore un'accidente in quanto in assenza della vista, come verrà specificato di seguito, non può esserci vero amore ma l'uomo esiste comunque.

Nella seconda quartina Giacomo sostiene infatti che sia sì possibile innamorarsi senza vedere l'oggetto del proprio amore ma che quell'amore non sia amore vero, critica in questo modo tutta una corrente letteraria e tematica detta \textit{amor de lonh}.

Nelle terzine torna ad avvalorare la sua tesi, dicendo in primo luogo che tutte le cose, belle o brutte, vengono prima filtrate dallo sguardo e per questo motivo l'esperienza amorosa è del tutto naturale, negli ultimi versi sostiene inoltre che sia universale il sentimento amoroso, andando così più in contro alla teoria dell'amore come sostanza
\end{quote}